Первые две главы дали язык: выпуклость, сильная выпуклость, липшицев градиент. Теперь --- сами методы. Общая схема $x_{k+1} = x_k + \alpha_k d_k$ распадается на два вопроса: куда шагать ($d_k$) и насколько ($\alpha_k$). Эта глава почти целиком про второй: сначала словарь (невязка, скорости сходимости), затем одномерная оптимизация --- точная (золотое сечение, параболы, Брент) и неточная (условия Армихо--Вульфа), а в конце первые теоремы о сходимости градиентного спуска с постоянным шагом $1/L$: сублинейная оценка для $C^{1,1}_L$ и линейная для сильно выпуклых функций.
